\section{Diccionario de los paquetes de diagnóstico}
Las fichas siguientes fijan el contenido, la frontera y la recepción de cada paquete. El Líder de Datos responde por su resultado; Análisis Funcional, Seguridad y Calidad aportan las tareas indicadas, con los titulares y funciones del SD1. El Líder de Datos ejerce la responsabilidad de migración prevista en SD5; su suplencia técnica se asignará entre los perfiles habilitados del equipo y se contabilizará en T-15. La contraparte funcional del cliente valida el significado de los antecedentes y las decisiones de negocio, sin asumir la ejecución técnica de la extracción. Administración designará su suplente funcional para reproducir el procedimiento durante el segundo ensayo completo, según SD5, sección 5.3.2.

\paragraph*{1.8.1.1 --- Catálogo de fuentes y accesos}
El entregable será un catálogo versionado de sistemas, archivos y expedientes con su custodio, ubicación, período, formato, cantidad de registros, tamaño, sensibilidad y mecanismo de acceso. Las cantidades que todavía no puedan comprobarse se identificarán como pendientes de extracción diagnóstica. El trabajo comprende preparar entrevistas, registrar fuentes y contrastar su cobertura con los seis conjuntos obligatorios; excluye digitalizar todos los documentos y sanear registros.

\textbf{Aceptación y responsable.} El Líder de Datos presentará una relación para cada conjunto, con sus fuentes identificadas o su ausencia documentada. Cada fuente tendrá custodio, acceso solicitado y condición de disponibilidad; las autorizaciones se conservarán como evidencia y las pendientes impedirán habilitar la extracción correspondiente. Administración, Escuela y Operaciones validarán los antecedentes de su ámbito.

\textbf{Esfuerzo y recursos.} Análisis Funcional: 16 horas para preparación y entrevistas; Datos: 8 horas para consolidación y contraste. Se requieren acceso autorizado al sistema de 2014, padrones, carpetas y registros escolares, y un repositorio protegido. 

\textbf{Supuestos, hitos y referencias.} Se estiman cuatro sesiones coordinadas: las dos iniciales de 90 minutos con quien conoce el legado, ya previstas en SD5, sección 5.3.2, y dos adicionales de hasta dos horas para contrastar las otras fuentes con sus custodios. Las 16 horas de Análisis Funcional incluyen hasta siete horas de sesiones y nueve de preparación y registro; no se supone dedicación exclusiva del cliente. El cierre mensual asistido previo al cambio administrativo se programará en los paquetes de conciliación, fuera de este diagnóstico. El catálogo identifica lotes y custodios, no revisa página por página los expedientes. Su recepción habilitará el diseño de extracción; depende de SUP-01 y SUP-02 del SD2 y de SD5, sección 5.3.1.

\paragraph*{1.8.1.2 --- Procedimiento de extracción diagnóstica}
El entregable será un procedimiento versionado con las herramientas necesarias para obtener una copia de diagnóstico y su manifiesto. Datos asignará 8 horas a examinar el mecanismo disponible, 16 a preparar la extracción y 8 a repetirla y documentarla; Seguridad asignará 8 horas a permisos, protección de copias y revisión del registro de acceso. El trabajo no modificará la fuente ni comprenderá cargas productivas.

\textbf{Aceptación y responsable.} El Líder de Datos demostrará dos ejecuciones repetibles sobre un origen o instantánea estable y autorizada. Los manifiestos permitirán comparar archivos o tablas, cantidades, períodos y sumas de comprobación; el procedimiento registrará versión de herramienta, parámetros, errores y ubicación protegida. Seguridad verificará que los permisos y las copias correspondan al alcance autorizado. Estas ejecuciones diagnósticas no sustituyen los dos ensayos completos de migración.

\textbf{Recursos.} Se requieren una copia protegida del origen, almacenamiento cifrado, herramientas de extracción y un técnico suplente que pueda reproducir el procedimiento. La repetición estará incluida en las 32 horas de Datos. Las licencias o insumos adicionales se identificarán por separado cuando el mecanismo elegido los requiera.

\textbf{Supuestos, hitos y referencias.} Las 40 horas suponen acceso a un mecanismo utilizable sin ingeniería inversa del ejecutable ni recuperación de soportes dañados. Esas condiciones se comprobarán al recibir el catálogo; su incumplimiento exigirá estimar los paquetes adicionales y revisar la programación. La disponibilidad del único operador conocedor del sistema se coordinará con un suplente de Synaptix. El procedimiento habilitará el perfilado y atenderá SD5, sección 5.3, y SUP-02 del SD2.

\paragraph*{1.8.1.3 --- Informe reproducible de calidad del origen}
El entregable será un informe acompañado de consultas o reglas versionadas para identificar campos obligatorios ausentes, duplicados, relaciones rotas, formatos incompatibles y diferencias entre registros y manifiestos. Datos dedicará 8 horas a reglas, 12 a ejecución y análisis y 4 a documentación; Calidad dedicará 8 horas a repetir comprobaciones y revisar la evidencia. El trabajo cuantificará incidencias sin corregir los registros de origen.

\textbf{Aceptación y responsable.} El Líder de Datos entregará resultados por conjunto con denominadores, cantidad de incidencias, referencia a registros afectados y limitaciones identificadas. Calidad reproducirá las comprobaciones sobre la misma copia y obtendrá los mismos resultados. Las diferencias monetarias se analizarán por moneda y período, sin compensar errores positivos y negativos. Los adjuntos se comprobarán por existencia y vínculo; su lectura semántica exhaustiva no pertenece a este paquete.

\textbf{Recursos.} Se requieren el catálogo, la copia de diagnóstico, sus manifiestos y un ambiente aislado de análisis. Las licencias o insumos adicionales se identificarán por separado cuando el mecanismo elegido los requiera.

\textbf{Supuestos, hitos y referencias.} Se supone que las fuentes tabulares pueden procesarse con las herramientas habilitadas en 1.8.1.2. La transcripción manual, OCR y reconstrucción documental se estimarán fuera de estas 32 horas. El informe habilitará las decisiones de saneamiento y conciliación; su referencia es SD5, secciones 5.2.4 y 5.3.2.

\paragraph*{1.8.1.4 --- Matriz de cobertura y decisiones de migración}
El entregable será una matriz que relacione cada conjunto obligatorio con sus fuentes, cobertura temporal, destino operativo o histórico, incidencias, reglas propuestas y validadores. Análisis Funcional dedicará 8 horas a decisiones y validación funcional; Datos, 8 a la matriz y sus vínculos; y Calidad, 8 a comprobar cobertura y frontera con los siguientes componentes.

\textbf{Aceptación y responsable.} El Líder de Datos entregará una fila por cada conjunto obligatorio y sus fuentes vinculadas, incluidos los 14 expedientes de abandono. Cada incidencia tendrá tratamiento propuesto y responsable de decisión. Los valores ausentes permanecerán identificados como tales; ninguna regla inventará pagos, autorizaciones ni movimientos históricos. Las decisiones aprobadas y las pendientes se distinguirán con su evidencia. La matriz se recibirá como diagnóstico completo cuando cubra todos los conjuntos; no autorizará una carga definitiva mientras existan decisiones bloqueantes sin resolver.

\textbf{Recursos.} Se requieren los tres entregables anteriores y la participación programada de los validadores del cliente. 

\textbf{Supuestos, hitos y referencias.} Las 24 horas comprenden una ronda consolidada de decisiones; los tiempos de espera del cliente se incorporarán al calendario. La matriz habilitará la especificación de saneamiento, carga y consulta histórica, con trazabilidad a SD3, sección 3.2, SD5, sección 5.3 y C07 RT-05.15.

Los hitos de recepción de los cuatro paquetes se enlazarán en el cronograma, manteniendo separados la entrega técnica de Synaptix y el tiempo de validación del cliente.
